\documentclass[11pt]{article}
\usepackage[letterpaper,margin=0.78in]{geometry}
\usepackage{amsmath,amssymb,mathtools}
\usepackage{array,booktabs,longtable}
\usepackage{microtype}
\usepackage{enumitem}
\usepackage[T1]{fontenc}
\usepackage{lmodern}
\usepackage{fancyhdr}
\usepackage{xcolor}
\usepackage{tcolorbox}
\usepackage{hyperref}
\hypersetup{hidelinks}
\setlength{\parindent}{0pt}
\setlength{\parskip}{6pt}
\setlist[itemize]{leftmargin=1.4em,itemsep=2pt,topsep=2pt}
\setlist[enumerate]{leftmargin=1.6em,itemsep=3pt,topsep=3pt}
\pagestyle{fancy}
\fancyhf{}
\lhead{LANGUAGE MECHANICS}
\rhead{CONDITION \& CONSTRAINT v0.2}
\cfoot{\thepage}
\newtcolorbox{lawbox}{colback=black!3,colframe=black!65,boxrule=0.45pt,arc=0pt,left=8pt,right=8pt,top=7pt,bottom=7pt}
\newcommand{\can}{\mathsf{can}}
\newcommand{\Refuse}{\varnothing}
\newcommand{\RR}{\mathbb{R}}

\begin{document}

\begin{center}
{\Large\bfseries CONDITION \& CONSTRAINT}\\[3pt]
{\large Structure Without Foundation}\\[8pt]
Anthony Vito Coppa\\[2pt]
\textit{Original proof composition: 22 January 2026}\\
\textit{Canonical standalone edition v0.2: 29 August 2026}
\end{center}

\begin{lawbox}
\textbf{Abstract.}
This paper isolates a minimum grammar of condition and constraint before truth, interpretation, quantity, or ontology is assigned. A condition declares an office in which a candidate can be addressed. A constraint declares what must continue to hold for that address to remain admissible under re-posing. Admission is local to a declared contract; non-vacuity requires at least one retained distinction; recognizability under reuse is typed by an equivalence relation at the declared resolution; and no downstream reader can reconstruct a distinction erased by an upstream chart. An inscriptional witness separates exact coincidence, nontrivial extension, and recognition without coincidence. A final conditional example proves the unique positive whole-remainder fixed point only after its additional chart contract is stated. The construction supplies no universal foundation. It supplies a discipline for saying where a rule applies, what it must preserve, and where continuation stops.
\end{lawbox}

\section{Position}

What follows is a constraint walk. Each step removes a form of statement that cannot survive its own declared use.

The paper does not ask what exists. It asks the prior and smaller question:

\begin{lawbox}
\centering
Under what declared conditions can a candidate be posed, retained, and posed again without the rule of admission silently changing?
\end{lawbox}

Existence, truth, reference, measurement, and meaning may become lawful downstream offices. None is required to draw the first boundary. The grammar therefore begins neither from a thing nor from a theory of things. It begins from a candidate address and a declared rule governing whether that address remains available.

\section{Condition and constraint}

\textbf{Definition 2.1 (Condition).}
A condition is a declared office of address
\[
D=(\Sigma,R,\rho),
\]
where:
\begin{itemize}
\item $\Sigma$ is a class of candidate determinations;
\item $R$ specifies the allowed relation of re-posing;
\item $\rho$ states the resolution at which retained distinctions are read.
\end{itemize}
No internal structure on $\Sigma$ is presumed.

\textbf{Definition 2.2 (Constraint).}
A constraint is a declared restriction on continuation in $D$. Equivalently, it specifies an admissible subclass
\[
A\subseteq\Sigma,
\]
together with whatever evidence the office requires for membership in $A$.

The condition fixes \emph{where} a term may lawfully apply. The constraint fixes \emph{what must continue to hold there}. Neither supplies content merely by being named.

Changing $\Sigma$, $R$, or $\rho$ creates a new version of the office. It does not retroactively alter the old one.

\textbf{Definition 2.3 (Admissibility).}
For $\sigma\in\Sigma$, write
\[
\can_D(\sigma)
\]
when $\sigma$ satisfies the declared constraint and remains available under the permitted re-posing relation.

When continuation fails, write $\Refuse$ at the metalevel. The mark $\Refuse$ is not an additional element of $\Sigma$.

Thus the first cut is binary only after the office has been fixed:
\[
\boxed{\text{continuation admitted}\quad|\quad\text{continuation refused}.}
\]
This is not a claim that the world is binary and it is not a valuation of the candidate. It is the yes-no output of one declared gate.

\section{Re-posing, recognition, and non-vacuity}

Let $R(\sigma,\tau)$ mean that $\tau$ is an allowed re-posing of $\sigma$ under $R$. The relation $R$ need not be temporal, causal, total, or deterministic. It may describe another inscription, another coordinate presentation, another run of a procedure, or another lawful use.

At the declared resolution $\rho$, let $\sim_\rho$ denote the equivalence relation induced by that resolution. Thus
\[
\sigma\sim_\rho\tau
\]
means that $\sigma$ and $\tau$ remain recognizable as belonging to the same declared class at resolution $\rho$.

This is not exact equality. It is not reference to an object outside the office. It is the minimum relation needed to say that re-posing has not become indistinguishable from replacement.

\textbf{Definition 3.1 (Non-vacuity).}
The office $D$ is non-vacuous at $\rho$ when there exist $\sigma,\tau\in\Sigma$ such that
\[
\can_D(\sigma),\qquad \can_D(\tau),\qquad R(\sigma,\tau),\qquad \sigma\not\sim_\rho\tau,
\]
while both addresses remain legible in the office.

This rules out the one-class quotient as a free victory. A constant gate certifies nothing about a distinction that the office promises to preserve. The floor is fixed by the promise, not loosened after the read.

\section{An inscriptional witness}

The abstract cut can be witnessed without introducing number or measure. Let $M$ be a class of legible marks and let $\oplus$ be a partial operation of re-inscription. Use:
\[
\equiv\ \text{for exact coincidence},\qquad
\neq\ \text{for non-coincidence},\qquad
\sim_\rho\ \text{for recognition at resolution }\rho.
\]

For a mark $x\in M$, exact closure may have the form
\[
x\oplus x\equiv x.
\]
Re-inscription has left no legible remainder at the declared resolution.

Nontrivial extension may have the form
\[
x\oplus x\neq x,
\qquad
x\sim_\rho(x\oplus x).
\]
The first relation retains difference; the second retains continuity of read. Together they witness extension without replacement.

\textbf{Proposition 4.1 (Witness triad).}
In this inscriptional chart, the three roles
\[
 x\equiv x,
 \qquad
 x\oplus x\neq x,
 \qquad
 x\sim_\rho(x\oplus x)
\]
are jointly sufficient to display coincidence, nontrivial reuse, and recognizability without coincidence.

\textit{Proof.}
The first expression supplies a stable coincidence test. The second supplies a pair distinguished at the declared resolution. The third prevents that pair from being read merely as unrelated replacement. No stronger classification is inferred. \hfill$\square$

The proposition is a witness, not an exhaustive ontology of persistence. It introduces no arithmetic law for $\oplus$, no identity element, and no claim that every admissible office has this presentation.

\section{The chart cannot restore what it erases}

A chart is not innocent merely because it is readable. Once a chart is used as the public carrier of an audit, its information loss bounds every later reader.

Let $H$ be a history class, let $\sim_\rho$ be the equivalence relation on $H$ induced by the declared resolution, and let
\[
\pi_\rho:H\to H/{\sim_\rho}
\]
be the quotient map. Let
\[
W:H\to\Sigma_W
\]
be a chart or writer. A reader is any function
\[
C:\Sigma_W\to H/{\sim_\rho}.
\]

\textbf{Lemma 5.1 (Erasure).}
If
\[
h_1\not\sim_\rho h_2
\qquad\text{but}\qquad
W(h_1)=W(h_2),
\]
then no reader $C$ can satisfy
\[
C\circ W=\pi_\rho
\]
on both $h_1$ and $h_2$.

\textit{Proof.}
Because $W(h_1)=W(h_2)$, functionality gives
\[
C(W(h_1))=C(W(h_2)).
\]
But $h_1\not\sim_\rho h_2$ implies
\[
\pi_\rho(h_1)\neq\pi_\rho(h_2).
\]
The two requirements are incompatible. \hfill$\square$

\begin{lawbox}
\centering
If a promised distinction is collapsed into one public mark, later interpretation cannot reconstruct it.
\end{lawbox}

The converse is equally important: preserving every raw difference is not automatically useful. The office must still declare which differences matter and how they are read. An admissible chart therefore needs both a non-degenerate resolution floor and a legible rule of return.

\section{Structure without foundation}

The word \emph{structure} names the relations that remain once a condition and its constraints have been declared. It does not name an absolute substrate beneath every possible office.

There is no required view from outside all conditions by which every condition must be ranked. That absence is a versioning rule rather than a missing hidden foundation.

Every assertion therefore carries its office:
\[
(\Sigma,R,\rho,A;\ \text{rule},\text{evidence},\text{result}).
\]

Two results may be compared only after a typed map between their offices has been supplied. Repeated vocabulary is not such a map.

The consequence is conservative:
\begin{itemize}
\item a result may be exact inside its declared condition without becoming universal;
\item a failed condition may be refused without converting that refusal into a claim about everything outside its domain;
\item a later construction may reuse a result only with its office and promised resolution intact.
\end{itemize}

\section{A conditional whole-remainder chart}

The preceding admissibility calculus does not select a special numerical coordinate. A numerical witness enters only after its own chart contract is stated.

Assume a positive coordinate $r$ is required to reproduce its own whole-remainder decomposition:
\[
r=1+\frac1r.
\]
This is a self-similar two-part chart. It is an additional condition, not a consequence of admissibility alone.

\textbf{Proposition 7.1 (Whole-remainder fixed point).}
Under the stated chart contract, the unique positive fixed point is
\[
\boxed{r=\frac{1+\sqrt5}{2}}.
\]

\textit{Proof.}
Multiplication by $r>0$ gives
\[
r^2-r-1=0.
\]
The quadratic formula gives
\[
r=\frac{1\pm\sqrt5}{2},
\]
of which exactly one root is positive. \hfill$\square$

Here $\sqrt5$ is the discriminant of this particular fixed-point equation. It does not follow from the bare admissibility cut, and this coordinate is not asserted to govern every non-closing traversal.

\section{Exact boundary}

The construction establishes only the following.

\begin{enumerate}
\item Admission is meaningful only relative to a declared condition.
\item A constraint states what must survive for continuation to remain admissible.
\item A non-vacuous office must preserve at least one distinction promised by its reader.
\item Reuse may retain recognizability without requiring exact coincidence.
\item No downstream reader can reconstruct a promised distinction erased by the chart that carries it.
\item Numerical or geometric witnesses enter only through separately stated chart contracts.
\item Comparison between offices requires an explicitly typed map; repeated words do not supply one.
\end{enumerate}

The paper establishes no theory of truth, semantics, ontology, dynamics, or universal geometry. It derives no physical law from admission.

The stopping point is exact:

\begin{lawbox}
\centering
A rule must say where it applies, what distinction it preserves, and what would count as failure.\\[3pt]
A reference may then orient a later use. Neither the rule nor the reference may supply missing content by implication.
\end{lawbox}

\section*{Provenance}

The formal argument was first composed on \textbf{22 January 2026}. This canonical standalone edition restates the argument as a self-contained Language Mechanics form. It introduces no external textual dependency and does not use later corpus results as premises.

\end{document}
